\subsection{The cubic bit in an order-216 marked affine quotient}

\paragraph{D23.P10 (a nine-fold smaller exact main gate).}
Let $V=\mathbb F_3^2$, $E_0=V\rtimes\operatorname{SL}_2(\mathbb F_3)$,
and use affine multiplication $(v,M)(w,N)=(v+Mw,MN)$.
There is a continuous surjection from the same marked pure free
group $F=\widehat F(x,y)$ as in E.65,
\[
 \xi:F\twoheadrightarrow E_0,\qquad
 \xi(x)=\left(0,\begin{pmatrix}-1&1\\0&-1\end{pmatrix}\right),\qquad
 \xi(y)=\left(e_1,\begin{pmatrix}-1&0\\-1&-1\end{pmatrix}\right).
\tag{D23.P10.1}
\]
The two linear matrices are $-U^2,-L^2$ in the source's positive
convention. For every actual multiplier-one gentle element
$h=(1,f)$ one has the exact formula
\[
 \boxed{\quad
 \xi(f)=\left(
   \begin{pmatrix}e\\\rho+(\rho+1)e\end{pmatrix},
   \begin{pmatrix}1&\rho\\\rho&1+\rho^2\end{pmatrix}
  \right),\qquad e=a_3(h)\in\{0,1\}\subset\mathbb F_3.
 \quad}
\tag{D23.P10.2}
\]
Thus the cubic bit is literally the first translation coordinate
of this order-$216$ finite evaluation of $f$.

The joint marked map
\[
 F\longrightarrow E_0\times S_4,\qquad
 x\longmapsto(\xi(x),(123)),\qquad
 y\longmapsto(\xi(y),(1234))
\tag{D23.P10.3}
\]
is onto, of order $5184$. Its image on $F'$ has order $864$ and
is exactly
\[
 (V\rtimes Q_8)\times A_4.
\tag{D23.P10.4}
\]
For actual full-GT elements, the supplied pentagon consequences
$a_f=\rho$ and $a_2=\nu_4$ leave precisely the same twelve
candidate states as D23.P4, now in a nine-fold smaller source
quotient. In these coordinates A.84 is exactly the assertion
that none of the six states with
\[
                  e\ne\nu_4
\tag{D23.P10.5}
\]
is attained by an actual full-GT element. The number of possible
mixed states has not decreased; the finite group realizing the
exact comparison has.

Among quotients of $E_0$, no proper quotient can retain this
detector on the supplied E.58 witness: every nontrivial normal
subgroup of $E_0$ contains $V$, whereas
\[
       \xi(f_{58})=((1,1)^{\mathsf T},I)\ne1.
\tag{D23.P10.6}
\]
This is a minimality statement within the quotient family of
$E_0$, not a claim of global minimality among all finite detectors.

\begin{proof}
First construct the map directly from the free presentation.
Write $X=\xi(x)$ and $Y=\xi(y)$. Their linear images generate
$\operatorname{SL}_2(\mathbb F_3)$: cubing either gives $-I$,
after which $U^2,L^2$ and hence $U,L$ are available. Moreover
\[
              [X^3,Y]=(e_1,I).
\]
Indeed $X^3=(0,-I)$, and conjugation by it negates the
translation of $Y$ while leaving its linear part unchanged;
the resulting commutator translation is $-2e_1=e_1$ in
characteristic three. Conjugating this translation by $Y$
gives the translation $(-1,-1)^{\mathsf T}$. Together these
span $V$. Thus the image is all of $E_0$, of order $9\cdot24=216$.

We next compare the marked map with D23.P3; this step fixes
the actual cubic sign rather than guessing from an isomorphic
finite group. On $E_F=\mathbb F_9^2\rtimes\operatorname{SL}_2(\mathbb F_3)$,
real-part projection of the translation is a group homomorphism,
because every linear matrix has entries in $\mathbb F_3$.
D23.P3's two pure generators have real translation $(1,1)$
and linear parts $-U^2,-L^2$. Change affine origin by
$t=(0,2)^{\mathsf T}$, which changes a translation $d$ of
linear part $M$ into $d-(I-M)t$. The two resulting translations
are $0,e_1$, precisely (D23.P10.1).

For an actual gentle $f$, D23.P3 gives
\[
 d_1=\frac{1-\rho-\epsilon}{1+i},\qquad
 d_2=(\rho+i)d_1,\qquad \epsilon=(-1)^{a_3(h)}.
\]
Since $(1+i)^{-1}=2+i$, writing
$c=1-\rho-\epsilon\in\mathbb F_3$ gives
\[
 \operatorname{Re}d_1=2c,\qquad
 \operatorname{Re}d_2=(2\rho-1)c.
\]
The origin correction is $(I-M_\rho)t=(\rho,\rho^2)^{\mathsf T}$.
If $(e,b)^{\mathsf T}$ denotes the corrected translation, then
\[
 e=2(1-\rho-\epsilon)-\rho,\qquad \epsilon=1+e.
\]
Thus $e=0$ when $a_3=0$ and $e=1$ when $a_3=1$.
Substituting $c=-\rho-e$ into the second coordinate gives
\[
 b=(2\rho-1)(-\rho-e)-\rho^2=\rho+(\rho+1)e.
\]
This proves (D23.P10.2), and E.58's $\rho=0,a_3=1$ gives
(D23.P10.6). The use of E.58 here is the supplied actual
witness; its absent infinite lifting input is not independently
reproved by this finite calculation.

Because $-I$ acts as $-1$ on $V$, commutators with it generate
all translations. Therefore $E_0'=V\rtimes Q_8$ and
$E_0^{\mathrm{ab}}\simeq C_3$. Since
$S_4^{\mathrm{ab}}\simeq C_2$ and both groups are solvable,
they have no nontrivial common quotient: such a quotient would
have a nontrivial abelian quotient common to $C_3,C_2$.
Goursat's lemma proves surjectivity of (D23.P10.3). Taking
derived subgroups proves (D23.P10.4). The further derived orders,
also verified in the exact certificate, are
\[
 |E_0^{(j)}|:216,72,18,9,1;\qquad
 |(E_0\times S_4)^{(j)}|:5184,864,72,9,1.
\]
On full GT, supplied FP.33 and FP.37 identify the marked
$S_4$ value as $t^{2\nu_4}p^\rho$ and identify $a_2=\nu_4$.
Together with (D23.P10.2), their equality criterion is exactly
(D23.P10.5). Reduction from arbitrary multiplier to multiplier
one is the supplied E.65.6 normalization, as in D23.P4.

Finally let $N\lhd E_0$ be nontrivial. If $N\cap V\ne1$,
the irreducibility of the natural $\operatorname{SL}_2(\mathbb F_3)$
module implies $V\subset N$. If $N\cap V=1$, then
$[N,V]\subset N\cap V$ forces $N$ to centralize $V$.
The centralizer of all translations in $E_0$ is $V$ itself,
since the linear action is faithful. Hence $N\subset V$,
contradicting $N\cap V=1$ and $N\ne1$. Every proper quotient
therefore kills the nonidentity translation in (D23.P10.6),
as claimed.

The exact program \path{agents/pentagon/cubic_detector_216.py}
checks the marked commutator, enumerates the affine image,
checks all six projected values from D23.P3, and independently
computes the order and derived series of the joint action on
nine affine points plus four permutation points.
\end{proof}
